\documentclass[a4paper]{article}
\usepackage{amsfonts}
\usepackage{amsmath}
\usepackage{amssymb}
\usepackage{graphicx}     
\newcommand{\dint}{\displaystyle\int}
\newcommand{\Bgtan}{\operatorname{Bgtan}}

\begin{document}
  
\noindent \large Auteur: Michael Livchits \\
\noindent \large Studierichting: Informatica\\
\noindent \large Oefeningenreeks 5A nr. 6.13\\

\medskip

\normalsize

$I = \displaystyle\dint_{0}^{\sqrt[4]{3}} \dfrac{2x}{1+x^2}dx$\\

Eerst vinden we de onbepaalde integraal\\

Stel $u = x^2$\\

$du = 2xdx$\\

$\displaystyle\int \dfrac{du}{1+u^2}$\\

$= \Bgtan u + C$\\

Pas grenzen aan en vul in\\

$I = \left[\Bgtan u\right]_{0}^{\sqrt{3}}$\\

$I = \Bgtan\sqrt{3} - \Bgtan0 = \dfrac{\pi}{3}$\\

\begin{thebibliography}{999}
%\bibitem{a} Referentie
\end{thebibliography}
\end{document} 